\documentclass[12pt,a4paper]{ctexart}
\usepackage[top=2.5cm,bottom=2.5cm,left=2.5cm,right=2.5cm]{geometry}
\usepackage{amsmath,amssymb,amsthm,physics}
\usepackage{graphicx,float,booktabs,array,caption}
\usepackage{hyperref,xcolor,enumitem}
\usepackage[numbers,sort&compress]{natbib}
\hypersetup{colorlinks=true,linkcolor=blue,citecolor=blue,urlcolor=blue}
\theoremstyle{definition}
\newtheorem{definition}{定义}[section]
\newtheorem{remark}{注记}[section]

\title{\heiti 格点QCD中的Symanzik有效理论：\\
从离散化误差的系统消除到连续极限的精度控制}
\author{基于本库文献的综合解析}
\date{\today}

\begin{document}
\maketitle

\begin{abstract}
\noindent
Symanzik有效理论是格点QCD中系统消除离散化误差的核心理论框架。其基本思想由Kurt Symanzik在1983年提出：格点作用量可被视为连续QCD作用量加上一系列由格点间距$a$的幂次控制的\textbf{无关算符}（irrelevant operators）的有效理论。通过在这些无关算符前添加适当选定的系数，可以在给定的$a$幂次上系统地消除离散化误差。本文系统阐述了Symanzik改进程序的完整理论体系：从连续理论中Symanzik有效作用量$\mathcal{L}_{\text{eff}} = \mathcal{L}_{\text{QCD}} + a\mathcal{L}^{(1)} + a^2\mathcal{L}^{(2)} + \cdots$的一般形式出发，详述规范部分的Symanzik改进——在标准Wilson plaquette作用量上添加$1\times2$矩形Wilson圈项，树级系数$c_0 = 5/3$, $c_1 = -1/12$消除$O(a^2)$误差；费米子部分的\textbf{Clover改进}（Sheikholeslami-Wohlert项）$\sim a c_{\text{SW}} \bar{\psi}\sigma_{\mu\nu}F_{\mu\nu}\psi$消除$O(a)$误差；\textbf{Tadpole改进}通过重标度$U_\mu \to U_\mu/u_0$重求微扰展开中的tadpole图贡献，改善微扰收敛性。重点展示CLQCD合作组在TITLS（Tadpole改进树级Symanzik规范作用量）和TITLC（Tadpole改进树级Clover费米子作用量）框架下的实践：从十一（或十三）个系综、五种格点间距的广泛计算，到$f_\pi$和$m_q$的线性$a^2$连续极限外推，再到$O(a\alpha_s)$残余修正和RI/MOM vs SMOM方案差异（$\sim7\%$）的系统处理。最后讨论Symanzik改进与HQET/LaMET中的重夸克/大动量展开之间的深刻类比对偶。
\end{abstract}

\tableofcontents
\newpage

\section{引言：格点上为何需要Symanzik改进？}

格点QCD将连续的Yang-Mills作用量和Dirac费米子作用量离散化为有限格点间距$a$上的格点作用量。在连续极限$a \to 0$下，格点理论还原为连续QCD。然而，在\textbf{有限格点间距}下，格点作用量与连续作用量之间存在系统性的差异——即\textbf{离散化误差}。

对于na\"ive的Wilson规范作用量\cite{CLQCD2025quark_mass,Chen2025gluon}：

\begin{equation}
S_W = \beta \sum_x \sum_{\mu<\nu} \left(1 - \frac{1}{N_c}\operatorname{Re}\operatorname{Tr} P_{\mu\nu}(x)\right), \qquad \beta = \frac{2N_c}{g^2},
\end{equation}

其离散化误差为$O(a^2)$。对于na\"ive的Wilson费米子作用量，离散化误差为$O(a)$（线性）。

\textbf{消除这些误差的意义：}$O(a)$的误差意味着将格点间距减半仅将误差减半——这是一个缓慢的收敛过程。$O(a^2)$的误差意味着格点间距减半将误差减少到原来的四分之一——显著更快。而$O(a^4)$（通过进一步的改进实现）则使收敛更快。对于当今的格点计算（$a \sim 0.05\text{--}0.15$~fm），$a^{-1} \sim 1.3\text{--}4$~GeV，未改进的$O(a)$误差在粗格点上的相对大小为$\mathcal{O}(a\Lambda_{\text{QCD}}) \sim 15\text{--}30\%$——这是一个不容忽视的系统误差。

\subsection{Symanzik改进的核心思想}

Kurt Symanzik于1983年提出的核心思想是：格点作用量可以被视为一个\textbf{有效理论}——连续QCD作用量加上一系列由格点间距$a$的幂次控制的\textbf{高维无关算符}：

\begin{definition}[Symanzik有效作用量]
\begin{equation}
\boxed{\mathcal{L}_{\text{eff}} = \mathcal{L}_{\text{QCD}} + a \mathcal{L}^{(1)} + a^2 \mathcal{L}^{(2)} + a^3 \mathcal{L}^{(3)} + \cdots},
\label{eq:Symanzik_effective_action}
\end{equation}
其中$\mathcal{L}^{(n)}$是由连续QCD的场和它们的导数构造的\textbf{量纲为$4+n$的局域算符}。通过在格点作用量中引入额外的项（具有适当调谐的系数），可以在给定的$a$幂次上系统地消除$\mathcal{L}^{(n)}$的贡献。
\end{definition}

\textbf{这一有效理论的严格性}来源于QCD的可重整化性质：在距离远大于格点间距的尺度上，所有``非物理''的格点效应可以被编码为局域的高维算符。随着$a \to 0$，这些算符的贡献被它们的量纲所决定的正幂次$a^n$压低。

\section{规范作用量的Symanzik改进}

\subsection{Wilson规范作用量的$O(a^2)$误差}

标准的Wilson规范作用量基于$1\times1$的plaquette。在连续极限展开中\cite{CLQCD2025quark_mass}：

\begin{equation}
S_W = \int d^4x\, \frac{1}{4}F^a_{\mu\nu}F^a_{\mu\nu} + a^2 \sum_x \sum_{\mu,\nu} c^{(2)}_{\mu\nu} \mathcal{O}^{(2)}_{\mu\nu}(x) + O(a^4),
\end{equation}

其中$\mathcal{O}^{(2)}$是量纲6的算符（规范场张量的四个导数组合）。$O(a^2)$项是主要误差来源。

\subsection{TILT规范作用量：树级Symanzik改进}

通过在作用量中引入\textbf{$1\times2$矩形Wilson圈}（长为$2a$，宽为$a$的矩形回路），可以在\textbf{树级}消除$O(a^2)$的离散化误差\cite{CLQCD2025quark_mass,Chen2025gluon}：

\begin{definition}[树级改进的Symanzik规范作用量]
\begin{equation}
\boxed{S_{\text{TILS}} = \beta \sum_x \left( c_0 \sum_{\mu<\nu} \left[1 - \frac{1}{3}\operatorname{Re}\operatorname{Tr} P_{\mu\nu}\right] + c_1 \sum_{\mu,\nu} \left[1 - \frac{1}{3}\operatorname{Re}\operatorname{Tr} R_{\mu\nu}\right] \right)},
\label{eq:TILS_action}
\end{equation}
其中$P_{\mu\nu}$是$1\times1$ plaquette，$R_{\mu\nu}$是$1\times2$矩形。\textbf{树级Symanzik系数}为：
\begin{equation}
\boxed{c_0 = \frac{5}{3}, \qquad c_1 = -\frac{1}{12}}。
\label{eq:tree_level_coefficients}
\end{equation}
\end{definition}

\textbf{系数推导的原理：}将$S_{\text{TILS}}$在$a \to 0$下展开为连续作用量加上$O(a^2)$修正。$c_0$和$c_1$被选择以使得$O(a^2)$修正的系数精确为零：

\begin{equation}
S_{\text{TILS}} = \int d^4x\, \frac{1}{4}F^a_{\mu\nu}F^a_{\mu\nu} + \underbrace{(c_0 + 8c_1 - 1)}_{=0\text{ 对树级系数}} \times O(a^2) + O(a^4)。
\end{equation}

\textbf{``树级''（tree-level）的含义：}这些系数仅在经典（无量子修正）水平上保证$O(a^2)$项的消除。量子修正（圈图效应）会产生额外的$O(a^2 g^2)$贡献，这些贡献可以通过非微扰调谐或tadpole改进来进一步压制。

\subsection{Tadpole改进：从TILS到TITLS}

\textbf{Tadpole改进}\cite{CLQCD2025quark_mass,CLQCD2024charm}是一种重求和微扰展开中tadpole图贡献的方法。在格点微扰论中，规范连接变量的高阶展开产生大量的tadpole图（一个胶子线从一个link出发并返回同一link），这些图产生大的$O(g^2)$修正，严重破坏了微扰展开的收敛性和树级Symanzik系数的有效性。

\textbf{Tadpole改进的核心操作：}

\begin{definition}[Tadpole改进的重标度]
\begin{equation}
\boxed{U_\mu(x) \to \frac{U_\mu(x)}{u_0}}, \qquad u_0 = \left\langle \frac{1}{3}\operatorname{Re}\operatorname{Tr} P_{\mu\nu} \right\rangle^{1/4}。
\label{eq:tadpole_rescaling}
\end{equation}
其中$u_0$是平均plaquette的\textbf{四次根}——这一选择的物理意义是：每个link变量$U_\mu$在规范不变的plaquette中贡献了四次，因此平均而言每个link的``大小''为$u_0$。
\end{definition}

\textbf{Tadpole改进的物理效应：}

\begin{enumerate}
    \item 所有规范连接变量被``缩小''了一个因子$u_0$（在典型格点间距下$u_0 \approx 0.8\text{--}0.9$），有效地吸收了tadpole图的主要贡献；
    \item 耦合常数被重定义为\textbf{tadpole改进的耦合}$\tilde{g}^2 = g^2/u_0^4$，在其展开下微扰级数的收敛性得到显著改善；
    \item CLQCD合作组使用的$\hat{\beta} = 10/(g^2 u_0^4)$正是tadpole改进后的有效规范耦合\cite{CLQCD2025quark_mass,CLQCD2024charm}。
\end{enumerate}

将tadpole改进应用于树级Symanzik作用量，就得到了\textbf{TITLS}（Tadpole Improved Tree-level Symanzik）规范作用量。

\section{费米子作用量的Symanzik改进：Clover项}

\subsection{Wilson费米子的$O(a)$问题}

na\"ive的Wilson费米子作用量具有$O(a)$的离散化误差——比规范部分的$O(a^2)$误差大一个量级。这是由Wilson项（为消除doubler费米子而添加的$ar\bar{\psi}\Box\psi$项，其中$r$是Wilson参数）引入的。

在连续极限下，Wilson项产生的$O(a)$修正可以编码为一个量纲5的有效算符：

\begin{equation}
\mathcal{L}^{(1)}_{\text{Wilson}} \propto a \, \bar{\psi} D_\mu D_\mu \psi。
\end{equation}

\subsection{Sheikholeslami-Wohlert（Clover）项}

Sheikholeslami和Wohlert于1985年证明：通过在Wilson费米子作用量中添加一个特定的量纲5算符——\textbf{Clover项}（或称Sheikholeslami-Wohlert项）——可以在树级消除\textbf{所有}$O(a)$的离散化误差：

\begin{definition}[Clover改进的费米子作用量]
\begin{equation}
\boxed{S_{\text{Clover}} = S_{\text{Wilson}} + a^5 \sum_x c_{\text{SW}} \, \bar{\psi}(x) \frac{i}{4}\sigma_{\mu\nu} F_{\mu\nu}(x) \psi(x)},
\label{eq:clover_fermion_action}
\end{equation}
其中$\sigma_{\mu\nu} = \frac{i}{2}[\gamma_\mu, \gamma_\nu]$，$F_{\mu\nu}$是由Clover plaquette组合构造的格点场强张量（量纲为$a^{-2}$），$c_{\text{SW}}$是\textbf{Sheikholeslami-Wohlert系数}。
\end{definition}

\textbf{物理图像：}Clover项描述了夸克自旋与胶子场强的Pauli型相互作用——类似于电子在电磁场中的$\vec{\sigma}\cdot\vec{B}$耦合。它通过在夸克传播中编码夸克-胶子相互作用的$O(a)$修正来抵消Wilson项引入的$O(a)$离散化误差。

\textbf{系数$c_{\text{SW}}$的值：}
\begin{itemize}
    \item \textbf{树级值：}$c_{\text{SW}} = 1$（在tadpole改进后）；
    \item \textbf{单圈改进值：}$c_{\text{SW}} = 1 + \mathcal{O}(g^2)$——通过格点微扰论计算单圈修正；
    \item \textbf{非微扰确定：}通过Schr\"odinger泛函（SF）方法或类似条件从实际格点数据中非微扰地调谐$c_{\text{SW}}$。
\end{itemize}

CLQCD合作组\cite{CLQCD2025quark_mass,CLQCD2024charm}使用Tadpole改进的树级Clover系数（TITLC），即$c_{\text{SW}}$取tadpole改进后的树级值。在连续极限$a \to 0$下，参数$c_{\text{SW}} \to 1$且$m_0 a \to 0$。

\subsection{轴矢流改进的省略：$O(a\alpha_s)$效应的回避}

CLQCD合作组\cite{CLQCD2025quark_mass}作出了一个有意识的选择：省略轴矢流的$O(a)$改进项：

\begin{quote}
``We have skipped the axial current improvement since the improvement coefficient itself strongly depends on the lattice spacing $a$ and bring the improvement term to be consistent with an $O(a^2)$ correction.''
\end{quote}

这一选择的物理动机在于：$O(a)$改进系数$c_A(g^2)$在格点间距$a$上具有强烈的依赖性。在CLQCD合作组使用的$a$值下，包含这一项实际上可能引入额外的$\sim O(a\alpha_s)$效应，而非简单的$O(a^2)$修正。``As evidence, both $f_\pi$ and $m_q$ show good consistency with a simple linear $a^2$ lattice spacing dependence''\cite{CLQCD2025quark_mass}——省略轴矢流改进后，物理观测量确实展现了干净的$a^2$标度行为。

\section{Tadple改进的深层结构：从微扰论到非微扰自洽性}

\subsection{平均plaquette $u_0$的物理意义}

Tadpole改进的核心量$u_0 = \langle \frac{1}{3}\operatorname{Re}\operatorname{Tr} P_{\mu\nu} \rangle^{1/4}$是一个\textbf{非微扰确定的量}——它通过对格点组态上所有plaquette取平均来计算。$u_0^4$度量了``量子涨落对经典plaquette值的压低程度''：

\begin{itemize}
    \item 在自由场极限（$g=0$，$U_\mu = 1$）下：$u_0 = 1$（无量子修正）；
    \item 在物理格点间距（$a \sim 0.1$~fm）下：$u_0 \approx 0.8\text{--}0.9$，反映了胶子涨落对Wilson圈显著的``面积定律''型压低。
\end{itemize}

\subsection{自洽tadpole改进：规范与费米子之间的对称性}

CLQCD合作组的TITLS+TITLC方案的一个关键优势是\textbf{自洽性}\cite{CLQCD2025quark_mass}：

\begin{quote}
``Both the fermion and gauge actions are tadpole improved self-consistently.''
\end{quote}

这意味着在规范作用量和费米子作用量中使用的是\textbf{相同的}tadpole因子$u_0$——从同一组态的plaquette中自洽确定。这保证了规范部分和费米子部分之间的理论一致性：如果两个部分使用不同的$u_0$值（或仅在其中一个中使用tadpole改进），则规范场和费米子场的真空涨落将在不同的有效格点间距上``运作''，导致额外的离散化效应。

\subsection{Tadpole改进耦合：$\hat{\beta}$}

CLQCD合作组使用tadpole改进的规范耦合常数为\cite{CLQCD2025quark_mass,CLQCD2024charm}：

\begin{definition}[Tadpole改进的规范耦合]
\begin{equation}
\boxed{\hat{\beta} = \frac{10}{g^2 u^4_0}}。
\label{eq:tadpole_beta}
\end{equation}
\end{definition}

不同的$\hat{\beta}$值对应于不同的格点间距系列：
\begin{itemize}
    \item \textbf{C系综：}$\hat{\beta}=6.20$，$a\approx0.105$~fm
    \item \textbf{F系综：}$\hat{\beta}=6.41$，$a\approx0.077$~fm
    \item \textbf{H系综：}$\hat{\beta}=6.72$，$a\approx0.052$~fm
\end{itemize}

\subsection{Tadpole改进与HQET/LaMET的类比}

一个深刻的洞察是\cite{CLQCD2024charm}：tadpole改进的逻辑与有效场论（HQET/LaMET）的逻辑高度一致。在HQET中，通过重定义重夸克场和改进正规化因子来吸收$m_Q a$修正\cite{CLQCD2024charm}：

\begin{quote}
``Similar to other effective theories, the convergence of the $O(m^2_Q a^2)$ corrections can be improved by redefining the expansion parameter $m_Q$ at a given cutoff of $1/a$, analogous to replacing $F_\pi(0)$ with $F_\pi(m_\pi)$ in chiral perturbation theory.''
\end{quote}

在tadpole改进中，通过重标度$U_\mu \to U_\mu/u_0$——类似于在手征微扰论中将$F_\pi(0)$替换为$F_\pi(m_\pi)$以改善展开的收敛性——来吸收微扰展开中大的tadpole修正。

\section{连续极限外推：Symanzik改进的实证检验}

\subsection{线性$a^2$外推的实证支持}

Symanzik改进的核心预期是：经过改进的格点理论中，主导的离散化误差应为$O(a^2)$（规范部分）和$O(a^2)$（Clover改进后的费米子部分，$O(a)$已被消除）。这意味着物理观测量$X$应在至少三个格点间距上被拟合为：

\begin{equation}
\boxed{X(a) = X(0) + d_X a^2 + \mathcal{O}(a^4)}。
\label{eq:a2_extrap}
\end{equation}

CLQCD合作组的数值结果\cite{CLQCD2025quark_mass}为这一预期提供了强有力的实证支持：
\begin{itemize}
    \item 在图3中，$(m_\pi)^2/m_q$和$f_\pi$在三个格点间距（$a = 0.105, 0.077, 0.052$~fm）上显示了\textbf{清晰的线性$a^2$行为}；
    \item $f_\pi$的格点间距依赖性较强（``continuum extrapolation pushes $f_\pi$ to be obviously higher''），而$(m^2_\pi/m_q)$比值对$a^2$的依赖性较弱——这是PCAC关系在格点上近似保持的体现。
\end{itemize}

\subsection{联合外推公式}

CLQCD合作组\cite{CLQCD2024charm}使用的完整联合外推公式：

\begin{equation}
X(m_\pi, m_{\eta_s}, a) = X(m^{\text{phys}}_\pi, m^{\text{phys}}_{\eta_s}, 0) + d^X_1 \Delta m^2_\pi + d^X_2 \Delta m^2_{\eta_s} + d^X_3 a^2 + d^X_4 a^4,
\label{eq:joint_extrap}
\end{equation}

同时包含了\textbf{物理夸克质量外推}（$\Delta m^2_\pi$, $\Delta m^2_{\eta_s}$项）和\textbf{连续极限外推}（$a^2$, $a^4$项）。

\subsection{残余$O(a\alpha_s)$效应的讨论}

尽管Clover改进消除了$O(a)$的误差，但\textbf{残余的$O(a\alpha_s)$交叉项}可能仍然存在，且在某些观测量中可能不可忽略\cite{CLQCD2025quark_mass}：

\begin{quote}
``Non-perturbative renormalization should remove all the $O(\alpha_s)$ effects, but not all the cross terms like the residual $O(a\alpha_s)$ effect of the clover action, which can cause the $O(a^2)$ continuum extrapolation to fail.''
\end{quote}

\textbf{解决方案：}CLQCD合作组利用RI/MOM与SMOM两个中间重整化方案之间的差异（在连续极限下约为$7\%$）作为残余离散化误差的系统估计——``Such a systematic uncertainty can also be considered as an estimate of the residual discretization error, as the correct continuum limit should be independent of the intermediate renormalization scheme.''\cite{CLQCD2025quark_mass}

\subsection{$O(a^4)$修正的必要性}

对于某些观测量——特别是\textbf{重夸克物理中的量}——$O(a^4)$项在联合外推中是必需的\cite{CLQCD2024charm}。这是因为：
\begin{itemize}
    \item 粲夸克质量$m_c \approx 1.3$~GeV在$a \approx 0.1$~fm处给出$(m_c a)^4 \approx 0.18$——$O(a^4 m^4_c)$修正不可忽略；
    \item $\overline{\text{MS}}$方案下的粲夸克质量对$O(a^4)$项特别敏感——需要通过包含$a^4$项或重夸克场的重定义来系统控制\cite{CLQCD2024charm}。
\end{itemize}

\section{Symanzik改进在CLQCD系综设计中的实践}

\subsection{系综的控制变量设计}

CLQCD合作组\cite{CLQCD2025quark_mass}的系综设计体现了Symanzik改进程序的系统化实践：

\begin{quote}
``The ensemble set used in this work is designed to control the variables in the systematic uncertainty estimation. For example, the spatial size $L$ of the C24P29, F32P30, and H48P32 ensembles are all within 1--2\% of each other, and the unitary pion masses are also similar with a 10\% difference. Thus, they are very suitable for investigating the discretization error of the hadron structure with non-zero given momentum.''
\end{quote}

\textbf{系综覆盖的参数空间}\cite{CLQCD2025quark_mass,CLQCD2024charm}：

\begin{table}[H]
\centering
\caption{CLQCD系综的Symanzik改进参数范围}
\label{tab:CLQCD_ensemble_range}
\begin{tabular}{lcc}
\toprule
参数 & 范围 & 控制目标 \\
\midrule
格点间距$a$       & $0.052\text{--}0.105$~fm（3--5个值） & 连续极限外推 \\
空间尺寸$L$       & $2.5\text{--}5.7$~fm           & 有限体积效应 \\
$\pi$介子质量     & $135\text{--}350$~MeV（7--8个值）& 手征外推 \\
$\eta_s$介子质量  & $644\text{--}749$~MeV           & 奇异夸克质量依赖性 \\
$m_\pi L$        & $2.6\text{--}5.7$                & 有限体积效应量化 \\
\bottomrule
\end{tabular}
\end{table}

\subsection{Clover系综的精度水平}

CLQCD合作组\cite{CLQCD2025quark_mass,CLQCD2024charm}的结果证明，经过Symanzik+Tadpole双重改进的Clover系综可以在广泛的物理观测量上实现与手征费米子（Domain Wall/Overlap）竞争甚至\textbf{优于}手征费米子的精度：

\begin{itemize}
    \item \textbf{轻夸克质量：}$m_u = 2.45(22)(20)$~MeV, $m_d = 4.74(11)(09)$~MeV, $m_s = 98.8(2.9)(4.7)$~MeV——与FLAG平均值在$1\text{--}2\sigma$内一致；
    \item \textbf{粲夸克质量：}$m_c(m_c) = 1.2933(72)(95)$~GeV——不确定度仅为$\sim 1\%$；
    \item \textbf{衰变常数：}$f_\pi = 130.7(0.9)(2.1)$~MeV, $f_K = 155.6(0.8)(2.7)$~MeV——精度已与当前最优的手征费米子结果相当。
\end{itemize}

\textbf{关键的结论：}``It is possible for the clover fermion to reach a high accuracy determination of the light quark masses''\cite{CLQCD2025quark_mass}。经过Symanzik改进和Tadpole改进双重处理的Clover费米子，可以以仅$\mathcal{O}(1/10\text{--}1/100)$ Domain Wall/Overlap费米子的计算代价，达到可比的物理精度。

\section{Symanzik改进、HQET与LaMET的有效场论统一视角}

\subsection{三种改进的对偶统一}

Symanzik改进、HQET和LaMET构成了一个有趣的\textbf{有效场论三元组}。它们都以``大标度展开''为核心方法，但在不同的标度和对象上进行操作：

\begin{table}[H]
\centering
\caption{三种有效场论改进的系统对偶}
\label{tab:three_EFT}
\begin{tabular}{lccc}
\toprule
要素 & Symanzik改进 & HQET & LaMET \\
\midrule
\textbf{大标度}     & $1/a$（UV截断） & $m_Q$（重夸克质量） & $P^z$（强子动量） \\
\textbf{消除对象}   & 离散化误差  & $m_Q$依赖性  & $P^z$依赖性 \\
\textbf{有效展开}   & $a^n \mathcal{L}^{(n)}$ & $1/m_Q^n$算符 & $1/(P^z)^n$幂次修正 \\
\textbf{匹配}       & 树级/单圈改进系数 & 微扰匹配$Z$ & 微扰匹配核$C$ \\
\textbf{``改进''操作}& 添加高维算符 & 重夸克场重定义 & Wilson线质量重正化 \\
\textbf{对标}       & \textit{连续QCD} & \textit{HQET} & \textit{光锥PDF} \\
\bottomrule
\end{tabular}
\end{table}

\subsection{共同的数学结构}

三种改进共享相同的\textbf{有效场论骨架}：

\begin{equation}
\mathcal{O}_{\text{全理论}} = Z(\text{大标度}) \otimes \mathcal{O}_{\text{有效理论}} + \mathcal{O}(\text{大标度}^{-\Delta})。
\end{equation}

\begin{itemize}
    \item \textbf{Symanzik：}格点观测量 = $Z(a\Lambda_{\text{QCD}}) \otimes$ 连续观测量 $+ \mathcal{O}(a^2)$；
    \item \textbf{HQET：}全QCD观测量 = $Z(m_Q/\Lambda_{\text{QCD}}) \otimes$ HQET观测量 $+ \mathcal{O}(1/m_Q)$\cite{Ji2014LaMET}；
    \item \textbf{LaMET：}准观测量 = $Z(P^z/\Lambda_{\text{QCD}}) \otimes$ 光锥观测量 $+ \mathcal{O}(1/(P^z)^2)$\cite{Ji2014LaMET}。
\end{itemize}

\textbf{Tadpole改进与HQET类比的深化：}在CLQCD合作组的粲物理计算\cite{CLQCD2024charm}中，重夸克改进的正规化因子$Z^c_V \equiv Z_V(\eta_c)$本质上是一种\textbf{tadpole改进在味空间的推广}——通过将矢量流正规化因子定义为``$\eta_c$介子态中的$Z_V$''来吸收有限格点间距下重夸克传播子中的$m_Q a$修正，完全类似于tadpole改进通过$u_0$吸收微扰展开中link变量的tadpole修正。

\section{总结与展望}

Symanzik有效理论是格点QCD中系统消除离散化误差的核心理论框架。本文的核心结论如下：

\begin{enumerate}
    \item \textbf{双层改进架构：}Symanzik改进在\textbf{规范部分}使用$c_0=5/3$, $c_1=-1/12$的$1\times2$矩形圈消除$O(a^2)$误差，在\textbf{费米子部分}使用$c_{\text{SW}}$的Clover项消除$O(a)$误差。二者结合后，主导的剩余误差降至$O(a^2)$量级\cite{CLQCD2025quark_mass}。

    \item \textbf{Tadpole改进是微扰收敛性的关键：}通过重标度$U_\mu \to U_\mu/u_0$和重定义$\hat{\beta} = 10/(g^2 u_0^4)$，tadpole改进吸收了微扰展开中最大的$O(g^2)$修正，使树级Symanzik系数在量子水平上仍然近似有效。自洽的Tadpole改进（规范+费米子）确保了理论与理论之间的对称性\cite{CLQCD2025quark_mass,CLQCD2024charm}。

    \item \textbf{实证验证：}CLQCD合作组在十三个系综、五种格点间距上的广泛计算提供了Symanzik+Tadpole双重改进的令人信服的实证验证。$f_\pi$和$m_q$展现了\textbf{清晰的线性$a^2$标度行为}，连续极限外推在$\chi^2/\text{d.o.f.} \sim 1$水平上与数据一致\cite{CLQCD2025quark_mass}。

    \item \textbf{残余误差的诚实估计：}$O(a\alpha_s)$交叉项的存在使得纯$a^2$外推可能不完全。CLQCD合作组通过RI/MOM vs SMOM方案差异（连续极限下$\sim 7\%$）来\textbf{honestly}估计残余离散化误差\cite{CLQCD2025quark_mass}。

    \item \textbf{有效场论的统一视角：}Symanzik改进、HQET和LaMET共享相同的有效场论骨架：大标度展开 + 微扰匹配 + 幂次修正。Tadpole改进与重夸克场重定义之间的类比揭示了三者之间深层的数学统一性\cite{Ji2014LaMET,CLQCD2024charm}。

    \item \textbf{经济性与精度的平衡：}Clover费米子在经过Symanzik+Tadpole双重改进后，可以以仅$\sim 1/10\text{--}1/100$手征费米子的计算代价达到可比的物理精度\cite{CLQCD2025quark_mass}——这对于需要在多格点间距上进行大量计算的PDF/胶子结构研究具有重大的实践意义。
\end{enumerate}

\textbf{未来方向：}更宽格点间距范围的系综以更好地约束残余$O(a\alpha_s)$效应；NNLO改进系数的非微扰确定（特别是对$c_{\text{SW}}$）；将Symanzik改进与梯度流正规化联合使用的系统探索。

\begin{thebibliography}{99}

\bibitem{CLQCD2025quark_mass}
H.-Y.~Du \textit{et al.} (CLQCD),
``Quark masses and low energy constants in the continuum from the tadpole improved clover ensembles,''
Phys.\ Rev.\ D \textbf{111}, 054504 (2025), arXiv:2408.03548.
\quad\textbf{[来源:} \texttt{docs/Quark masses and low energy constants in the continuum from the tadpole improved clover ensembles.pdf}\textbf{]}
（含TITLS+TITLC作用量的详细构造、$O(a\alpha_s)$残余误差讨论、11系综联合外推）

\bibitem{CLQCD2024charm}
H.-Y.~Du \textit{et al.} (CLQCD),
``Charmed meson masses and decay constants in the continuum from the tadpole improved clover ensembles,''
Phys.\ Rev.\ D \textbf{109}, 054507 (2024), arXiv:2310.00814.
\quad\textbf{[来源:} \texttt{docs/Charmed meson masses and decay constants in the continuum from the tadpole improved clover ensembles.pdf}\textbf{]}
（含重夸克改进正规化与Symanzik改进的有效场论类比、13系综联合外推）

\bibitem{Chen2025gluon}
C.~Chen \textit{et al.} (CLQCD, Lattice Parton),
``Unpolarized gluon PDF of the nucleon from lattice QCD in the continuum limit,''
(2025), arXiv:2510.26425.
\quad\textbf{[来源:} \texttt{docs/Unpolarized gluon PDF of the nucleon from lattice QCD in the continuum limit.pdf}\textbf{]}
（含TITLS+TITLC作用量在胶子PDF连续极限计算中的应用）

\bibitem{Chen2025first}
C.~Chen \textit{et al.} (CLQCD, Lattice Parton),
``First self-renormalized gluon PDF of nucleon from large-momentum effective theory in the continuum limit,''
Phys.\ Rev.\ D \textbf{111}, 074506 (2025), arXiv:2408.12819.
\quad\textbf{[来源:} \texttt{docs/First Self-Renormalized Gluon PDF of Nucleon from Large-Momentum Effective Theory in the Continuum Limit.pdf}\textbf{]}

\bibitem{Ji2014LaMET}
X.~Ji,
``Parton physics from large-momentum effective field theory,''
Sci.\ China Phys.\ Mech.\ Astron.\ \textbf{57}, 1407 (2014), arXiv:1404.6680.
\quad\textbf{[来源:} \texttt{docs/Parton Physics from Large-Momentum Effective Field Theory.pdf}\textbf{]}
（含HQET与LaMET的类比——有效场论框架的统一视角）

\bibitem{Huo2021self}
Y.-K.~Huo \textit{et al.} (LPC),
``Self-renormalization of quasi-light-front correlators on the lattice,''
Nucl.\ Phys.\ B \textbf{969}, 115443 (2021), arXiv:2103.02965.
\quad\textbf{[来源:} \texttt{docs/Self-Renormalization of Quasi-Light-Front Correlators on the Lattice.pdf}\textbf{]}
（含Clover与Overlap价夸克在不同格点作用量上的跨检查）

\bibitem{Zhang2024gauge_fixing}
K.~Zhang \textit{et al.} (LPC),
``Impact of gauge fixing precision on the continuum limit of non-local quark-bilinear lattice operators,''
Phys.\ Rev.\ D \textbf{110}, 074505 (2024), arXiv:2405.14097.
\quad\textbf{[来源:} \texttt{docs/Impact of gauge fixing precision on the continuum limit of non-local quark-bilinear lattice operators.pdf}\textbf{]}
（含5种格点间距上规范固定精度对连续极限的系统研究）

\bibitem{NieMiera2025}
A.~NieMiera, W.~Good, H.-W.~Lin, and F.~Yao,
``First self-renormalized gluon PDF of nucleon from large-momentum effective theory in the continuum limit,''
(2025), arXiv:2510.17758.
\quad\textbf{[来源:} \texttt{docs/First Self-Renormalized Gluon PDF of Nucleon from Large-Momentum Effective Theory in the Continuum Limit.pdf}\textbf{]}

\bibitem{Egerer2022}
C.~Egerer \textit{et al.} (HadStruc),
``Towards the determination of the gluon helicity distribution,''
Phys.\ Rev.\ D \textbf{106}, 094511 (2022), arXiv:2207.08733.
\quad\textbf{[来源:} \texttt{docs/Towards the determination of the gluon helicity distribution in the nucleon from lattice quantum chromodynamics.pdf}\textbf{]}

\bibitem{NoteGluonPDFs}
``Note of gluon PDFs,'' 内部笔记。
\quad\textbf{[来源:} \texttt{docs/Note of gluon PDFs.pdf}\textbf{]}

\bibitem{Fan2018gluon}
Z.-Y.~Fan, Y.-B.~Yang, A.~Anthony, H.-W.~Lin, and K.-F.~Liu,
``Gluon quasi-PDF from lattice QCD,''
Phys.\ Rev.\ Lett.\ \textbf{121}, 242001 (2018), arXiv:1808.02077.
\quad\textbf{[来源:} \texttt{docs/Gluon Quasi-PDF From Lattice QCD.pdf}\textbf{]}

\bibitem{Ji2013Euclidean}
X.~Ji,
``Parton physics on a Euclidean lattice,''
Phys.\ Rev.\ Lett.\ \textbf{110}, 262002 (2013), arXiv:1305.1539.
\quad\textbf{[来源:} \texttt{docs/Parton Physics on Euclidean Lattice.pdf}\textbf{]}

\bibitem{Good2025a}
W.~Good, F.~Yao, and H.-W.~Lin,
``First nucleon gluon PDF from large momentum effective theory,''
(2025), arXiv:2505.13321.
\quad\textbf{[来源:} \texttt{docs/First Nucleon Gluon PDF from Large Momentum Effective Theory.pdf}\textbf{]}

\bibitem{Lin2025gluon}
W.~Good, K.~Hasan, and H.-W.~Lin,
``Gluon parton distribution of the nucleon from $2+1+1$-flavor lattice QCD in the physical-continuum limit,''
J.\ Phys.\ G \textbf{52}, 035105 (2025), arXiv:2409.02750.
\quad\textbf{[来源:} \texttt{docs/Gluon Parton Distribution of the Nucleon from 2+1+1-Flavor Lattice QCD in the Physical-Continuum Limit.pdf}\textbf{]}

\end{thebibliography}

\end{document}
